% Öffentliche Aussagen bleiben in Band 37; hier werden sie unverändert wiederholt.
\section*{Index und Periode endlicher Halbgruppen}
\hypertarget{index-period.proofs}{}
Diese Ergänzung führt von der ersten Wiederholung einer positiven Potenzfolge
zum eindeutigen Index--Perioden-Paar, zur periodischen Fortsetzung und zu
idempotenten Potenzen. Die \IndexPeriodReadingLink{} erklärt die Argumente
an Beispielen; die \IndexPeriodMainLink{} enthält die ursprünglichen
Definitionen, Axiome und Sätze mit unveränderten Kennungen und Nummern.
Vor jedem Beweis steht seine vollständige Aussage. Die hier wiederholten
öffentlichen Aussagen werden nicht erneut registriert. Zusätzliche
Wohldefiniertheitssätze tragen die Bandkennung \(37\mathrm I\).

Alle Potenzen haben positive Exponenten. Im Fortsetzungsargument wird
\(t=0\) gesondert behandelt; eine Potenz \(a^0\) wird nicht verwendet.
Die Aussagen über ein Element setzen \(a\in A\) voraus und sind für den
leeren Träger daher ohne Anwendungsfall. Erst die abschließende Existenz
eines Idempotenten verlangt ausdrücklich \(A\neq\varnothing\).

\Needspace{24\baselineskip}
\section{Wiederholungsstellen und das eindeutige Parameterpaar}
\subsection{Wohldefiniertheit der Wiederholungsmenge}
\hypertarget{index-period.proof.FiniteSemigroupPowerRepetitionSetDef}{}
\IndexPeriodStatement[Definition]{FiniteSemigroupPowerRepetitionSetDef}
\Needspace{16\baselineskip}
\FormulaThmDeltaK[Die Potenzen im Wiederholungsprädikat sind typisiert]{%
 \begin{aligned}[t]
 &\SemigroupAlg{A}{\star}\dsep a\in A\dsep i,j\in\mathbb N
 \dsep i\neq0\dsep j\neq0\vdash{}\\[-2pt]
 &a^i\in A\land a^j\in A
 \end{aligned}
}{B37IPPositivePowersTyped}{\DeltaRow{Mengen}{A\dsep a\dsep i\dsep j}\DeltaRow{Funktionen}{\star}}
\begin{tabproofwide}
 \proofstepwidestar[1]{\SemigroupAlg{A}{\star}}{\rA}
 \proofstepwidestar[2]{a\in A}{\rA}
 \proofstepwidestar[3]{i,j\in\mathbb N}{\rA}
 \proofstepwidestar[4]{i\neq0}{\rA}
 \proofstepwidestar[5]{j\neq0}{\rA}
 \proofstepwidestar[1,2,3,4]{a^i\in A}{\FormulaRefAuto{SemigroupPowerClosure}[thm]{1,2,\rAEa{3},4}}
 \proofstepwidestar[1,2,3,5]{a^j\in A}{\FormulaRefAuto{SemigroupPowerClosure}[thm]{1,2,\rAEb{3},5}}
 \proofstepwidestar[1,2,3,4,5]{a^i\in A\land a^j\in A}{\rAI{6,7}}
\end{tabproofwide}
Im Wiederholungsprädikat werden beide Potenzen unter genau diesen
Typ- und Positivitätsbedingungen verglichen. Eine Potenz mit Exponent null
ist somit für die Definition nicht erforderlich.

\Needspace{18\baselineskip}
\FormulaThmDeltaK[Die Wiederholungsmenge ist eindeutig bestimmt]{%
 \begin{aligned}[t]
 &\SemigroupAlg{A}{\star}\dsep a\in A\vdash{}\\[-2pt]
 &\Bigl(\exists!R\,\forall j\,\Bigl(j\in R\leftrightarrow
     \bigl(j\in\mathbb N\land j\neq0\land{}\\[-2pt]
 &\qquad\exists i\in\mathbb N\,
       (i\neq0\land i<j\land a^i=a^j)\bigr)\Bigr)\Bigr)
       \land\mathcal R_{A,\star}(a)\subseteq\mathbb N
 \end{aligned}
}{B37IPRepetitionSetWellDefined}{\DeltaRow{Mengen}{A\dsep a\dsep i\dsep j\dsep R}\DeltaRow{Funktionen}{\star}}
In dieser Tabelle kürzen wir
\[
 Q(j)\equiv j\neq0\land\exists i\in\mathbb N\,
   (i\neq0\land i<j\land a^i=a^j)
\]
ab. Die eindeutige Existenz folgt aus Aussonderung und Extensionalität;
sie benötigt noch keine Endlichkeit von \(A\).
\begin{tabproofwide}
 \proofstepwidestar[1]{\SemigroupAlg{A}{\star}}{\rA}
 \proofstepwidestar[2]{a\in A}{\rA}
 \proofstepwidestar[]{\exists R\,\forall j\,(j\in R\leftrightarrow j\in\mathbb N\land Q(j))}{\FormulaRefAuto{\exists B\;\forall x\;\bigl(x \in B \;\leftrightarrow\; x \in A \,\land\, P(x)\bigr)}[ax]}
 \proofstepwidestar[]{\exists!R\,\forall j\,(j\in R\leftrightarrow j\in\mathbb N\land Q(j))}{\FormulaRefAuto{\exists A(\forall x(x \in A\leftrightarrow P(x)))\vdash \exists! A(\forall x(x \in A\leftrightarrow P(x)))}[thm]{3}}
 \proofstepwidestar[5]{j\in\mathcal R_{A,\star}(a)}{\rA}
 \proofstepwidestar[1,2,5]{j\in\{v\in\mathbb N\mid Q(v)\}}{\FormulaRefAuto{FiniteSemigroupPowerRepetitionSetDef}[def]{1,2,5}}
 \proofstepwidestar[1,2,5]{j\in\mathbb N\land Q(j)}{\FormulaRefAuto{x\in\{u\in A\mid P(u)\}\eqvdash x\in A\land P(x)}[thm]{6}}
 \proofstepwidestar[1,2,5]{j\in\mathbb N}{\rAEa{7}}
 \proofstepwidestar[1,2]{\mathcal R_{A,\star}(a)\subseteq\mathbb N}{\FormulaRefAuto{A\subseteq B\coloneq\forall x\,(x\in A\to x\in B)}[def]{\rUI{\rRI{5,8}}}}
 \proofstepwidestar[1,2]{\begin{aligned}[t]&\bigl(\exists!R\,\forall j\,(j\in R\leftrightarrow j\in\mathbb N\land Q(j))\bigr)\\[-2pt]&{}\land\mathcal R_{A,\star}(a)\subseteq\mathbb N\end{aligned}}{\rAI{4,9}}
\end{tabproofwide}


\subsection{Eine endliche Potenzfolge wiederholt sich}
\hypertarget{index-period.proof.FiniteSemigroupPowerRepetitionExists}{}
\IndexPeriodStatement{FiniteSemigroupPowerRepetitionExists}
Wir schreiben \(\sigma=\operatorname{pow}_{A,\star}(a)\) für die Potenzfolge aus Band~28.
Ihre Werte erfüllen \(\sigma(k)=a^{k+1}\). Die beiden benötigten
Auswertungen werden in der Tabelle einzeln eingesetzt.
\begin{tabproofwide}
 \proofstepwidestar[1]{\FiniteSemigroupAlg{A}{\star}}{\rA}
 \proofstepwidestar[2]{a\in A}{\rA}
 \proofstepwidestar[1]{\SemigroupAlg{A}{\star}}{\FormulaRefAuto{FiniteSemigroupBaseAxiom}[ax]{1}}
 \proofstepwidestar[1]{\Finite A}{\FormulaRefAuto{FiniteSemigroupFinitenessAxiom}[ax]{1}}
 \proofstepwidestar[1,2]{\sigma\colon\mathbb N\to A}{\FormulaRefAuto{SemigroupPowerSequenceFunction}[thm]{3,2}}
 \proofstepwidestar[1,2]{\exists k\,\exists n\,(k\in\mathbb N\land n\in\mathbb N\land k<n\land\sigma(k)=\sigma(n))}{\FormulaRefAuto{FiniteSequenceRepetition}[thm]{4,5}}
 \proofstepwidestar[7]{\exists n\,(k\in\mathbb N\land n\in\mathbb N\land k<n\land\sigma(k)=\sigma(n))}{\rA}
 \proofstepwidestar[8]{k\in\mathbb N\land n\in\mathbb N\land k<n\land\sigma(k)=\sigma(n)}{\rA}
 \proofstepwidestar[]{1\in\mathbb N}{\FormulaRefAuto{PeanoOneInN}[thm]}
 \proofstepwidestar[8]{k\in\mathbb N}{\rAEa{8}}
 \proofstepwidestar[8]{n\in\mathbb N}{\rAEa{\rAEb{8}}}
 \proofstepwidestar[1,2,8]{\sigma(k)=a^{k+1}}{\FormulaRefAuto{a=b\vdash b=a}[thm]{\FormulaRefAuto{SemigroupPowerSuccessorIndex}[thm]{3,2,10}}}
 \proofstepwidestar[1,2,8]{\sigma(n)=a^{n+1}}{\FormulaRefAuto{a=b\vdash b=a}[thm]{\FormulaRefAuto{SemigroupPowerSuccessorIndex}[thm]{3,2,11}}}
 \proofstepwidestar[1,2,8]{a^{k+1}=a^{n+1}}{\rIE{12,13,\rAEb{\rAEb{\rAEb{8}}}}}
 \proofstepwidestar[8]{k+1\in\mathbb N}{\FormulaRefAuto{PeanoAddClosure}[thm]{10,9}}
 \proofstepwidestar[8]{n+1\in\mathbb N}{\FormulaRefAuto{PeanoAddClosure}[thm]{11,9}}
 \proofstepwidestar[8]{k+1\neq0}{\FormulaRefAuto{PeanoAddOneNonzero}[thm]{10}}
 \proofstepwidestar[8]{n+1\neq0}{\FormulaRefAuto{PeanoAddOneNonzero}[thm]{11}}
 \proofstepwidestar[8]{1+k<1+n}{\FormulaRefAuto{PeanoAddLeftStrictMonotone}[thm]{9,10,11,\rAEa{\rAEb{\rAEb{8}}}}}
 \proofstepwidestar[8]{k+1=1+k}{\FormulaRefAuto{PeanoAddCommutative}[thm]{10,9}}
 \proofstepwidestar[8]{n+1=1+n}{\FormulaRefAuto{PeanoAddCommutative}[thm]{11,9}}
 \proofstepwidestar[8]{k+1<n+1}{\rIE{\FormulaRefAuto{a=b\vdash b=a}[thm]{20},\FormulaRefAuto{a=b\vdash b=a}[thm]{21},19}}
 \proofstepwidestar[1,2,8]{\exists i\in\mathbb N\,(i\neq0\land i<n+1\land a^i=a^{n+1})}{\rEI{\rAI{15,\rAI{17,\rAI{22,14}}}}}
 \proofstepwidestar[1,2,8]{\exists j\in\mathbb N\,(j\neq0\land\exists i\in\mathbb N\,(i\neq0\land i<j\land a^i=a^j))}{\rEI{\rAI{16,\rAI{18,23}}}}
 \proofstepwidestar[1,2,7]{\exists j\in\mathbb N\,(j\neq0\land\exists i\in\mathbb N\,(i\neq0\land i<j\land a^i=a^j))}{\rEE{7,8,24}}
 \proofstepwidestar[1,2]{\exists j\in\mathbb N\,(j\neq0\land\exists i\in\mathbb N\,(i\neq0\land i<j\land a^i=a^j))}{\rEE{6,7,25}}
 \proofstepwidestar[27]{j\in\mathbb N\land(j\neq0\land\exists i\in\mathbb N\,(i\neq0\land i<j\land a^i=a^j))}{\rA}
 \proofstepwidestar[27]{\exists i\in\mathbb N\,(i\neq0\land i<j\land a^i=a^j)}{\rAEb{\rAEb{27}}}
 \proofstepwidestar[29]{i\in\mathbb N\land(i\neq0\land i<j\land a^i=a^j)}{\rA}
 \proofstepwidestar[1,2,27,29]{j\in\mathcal R_{A,\star}(a)}{\FormulaRefAuto{FiniteSemigroupPowerRepetitionSetDef}[def]{3,2,\rAEa{27},\rAEa{\rAEb{27}},\rEI{29}}}
 \proofstepwidestar[1,2,27,29]{\mathcal R_{A,\star}(a)\neq\varnothing}{\FormulaRefAuto{x\in A\vdash A\neq\varnothing}[thm]{30}}
 \proofstepwidestar[1,2,27]{\mathcal R_{A,\star}(a)\neq\varnothing}{\rEE{28,29,31}}
 \proofstepwidestar[1,2]{\mathcal R_{A,\star}(a)\neq\varnothing}{\rEE{26,27,32}}
\end{tabproofwide}


\Needspace{24\baselineskip}
\subsection{Die Eigenschaften des Index--Perioden-Paares}
Der Begriff ist durch die folgenden drei Bedingungen festgelegt.
Diese Bedingungen sind die Zugriffsaxiome des Begriffs; sie werden nicht
als zusätzliche Existenzbehauptungen verwendet. Die Existenz eines solchen
Paares und seine Eindeutigkeit werden anschließend bewiesen.
\IndexPeriodStatement[Definition]{FiniteSemigroupIndexPeriodPairDef}
\IndexPeriodStatement[Voraussetzungen des Begriffs]{FiniteSemigroupIndexPeriodPositiveData}
\IndexPeriodStatement[Voraussetzungen des Begriffs]{FiniteSemigroupIndexPeriodEquality}
\IndexPeriodStatement[Voraussetzungen des Begriffs]{FiniteSemigroupIndexPeriodInitialDistinctness}

\hypertarget{index-period.proof.FiniteSemigroupLeastRepetitionInitialDistinctness}{}
\IndexPeriodStatement{FiniteSemigroupLeastRepetitionInitialDistinctness}
% Aus Band 37 ausgelagerte Beweistabelle.
\begin{tabproofwide}
  \proofstepwidestar[1]{\SemigroupAlg{A}{\star}}{\rA}
  \proofstepwidestar[2]{a\in A}{\rA}
  \proofstepwidestar[3]{j\in\mathbb N}{\rA}
  \proofstepwidestar[4]{%
    \forall v\in\mathcal R_{A,\star}(a)\,(j\leq v)}{\rA}
  \proofstepwidestar[5]{r,s\in\mathbb N}{\rA}
  \proofstepwidestar[6]{r\neq0}{\rA}
  \proofstepwidestar[7]{r<s}{\rA}
  \proofstepwidestar[8]{s<j}{\rA}
  \proofstepwidestar[9]{a^r=a^s}{\rA}
  \proofstepwidestarbreak[5,7]{s\neq0}{%
    \FormulaRefAuto{PeanoStrictUpperNonzero}[thm]{%
      \rAEa{5},\rAEb{5},7}}
  \proofstepwidestarbreak[1,2,5,6,7,9]{%
    s\in\mathcal R_{A,\star}(a)}{%
    \FormulaRefAuto{FiniteSemigroupPowerRepetitionSetDef}[def]{%
      1,2,\rAEb{5},10,
      \rEI{\rAI{\rAEa{5},\rAI{6,\rAI{7,9}}}}}}
  \proofstepwidestarbreak[1,2,4,5,6,7,9]{j\leq s}{%
    \rRE{11,\rUE{4}}}
  \proofstepwidestarbreak[1,2,3,4,5,6,7,8,9]{\bot}{%
    \FormulaRefAuto{PeanoOrderNoCycle}[thm]{\rAEb{5},3,8,12}}
  \proofstepwidestar[1,2,3,4,5,6,7,8]{a^r\neq a^s}{\rCI{9,13}}
\end{tabproofwide}


\hypertarget{index-period.proof.FiniteSemigroupLeastRepetitionProducesIndexPeriod}{}
\IndexPeriodStatement{FiniteSemigroupLeastRepetitionProducesIndexPeriod}
% Aus Band 37 ausgelagerte Beweistabelle.
\begin{tabproofwide}
  \proofstepwidestar[1]{\SemigroupAlg{A}{\star}}{\rA}
  \proofstepwidestar[2]{a\in A}{\rA}
  \proofstepwidestar[3]{i,j,q\in\mathbb N}{\rA}
  \proofstepwidestar[4]{i\neq0}{\rA}
  \proofstepwidestar[5]{j\neq0}{\rA}
  \proofstepwidestar[6]{i<j}{\rA}
  \proofstepwidestar[7]{a^i=a^j}{\rA}
  \proofstepwidestar[8]{%
    \forall v\in\mathcal R_{A,\star}(a)\,(j\leq v)}{\rA}
  \proofstepwidestar[9]{i+(q+1)=j}{\rA}
  \proofstepwidestar[]{1\in\mathbb N}{%
      \FormulaRefAuto{PeanoOneInN}[thm]}
  \proofstepwidestarbreak[3]{q+1\in\mathbb N}{%
    \FormulaRefAuto{PeanoAddClosure}[thm]{%
      \rAEb{\rAEb{3}},10}}
  \proofstepwidestarbreak[3]{q+1\neq0}{%
    \FormulaRefAuto{PeanoAddOneNonzero}[thm]{\rAEb{\rAEb{3}}}}
  \proofstepwidestarbreak[7,9]{a^{i+(q+1)}=a^i}{%
    \rChain{%
      \rIE{9,\rII},
      \FormulaRefAuto{a=b\vdash b=a}[thm]{7}}}
  \proofstepwidestar[14]{r\in\mathbb N}{\rA}
  \proofstepwidestar[15]{s\in\mathbb N}{\rA}
  \proofstepwidestar[16]{%
    r\neq0\land r<s\land s<i+(q+1)}{\rA}
  \proofstepwidestarbreak[9,16]{s<j}{%
    \rIE{9,\rAEb{\rAEb{16}}}}
  \proofstepwidestarbreak[1,2,3,8,9,14,15,16]{a^r\neq a^s}{%
    \FormulaRefAuto{FiniteSemigroupLeastRepetitionInitialDistinctness}[thm]{%
      1,2,\rAEa{\rAEb{3}},8,\rAI{14,15},
      \rAEa{16},\rAEa{\rAEb{16}},17}}
  \proofstepwidestarbreak[1,2,3,8,9]{%
    \forall r,s\in\mathbb N\,
      \bigl(r\neq0\land r<s\land s<i+(q+1)
        \rightarrow a^r\neq a^s\bigr)}{%
    \rUI{\rRI{14,\rUI{\rRI{15,\rRI{16,18}}}}}}
  \proofstepwidestarbreak[1,2,3,4,7,8,9]{%
    \operatorname{IP}_{A,\star}(a;i,q+1)}{%
    \FormulaRefAuto{FiniteSemigroupIndexPeriodPairDef}[def]{%
      1,2,\rAEa{3},4,11,12,13,19}}
\end{tabproofwide}


\hypertarget{index-period.proof.FiniteSemigroupIndexPeriodExistence}{}
\IndexPeriodStatement{FiniteSemigroupIndexPeriodExistence}
% Aus Band 37 ausgelagerte Beweistabelle.
\begin{tabproofwide}
  \proofstepwidestar[1]{\FiniteSemigroupAlg{A}{\star}}{\rA}
  \proofstepwidestar[2]{a\in A}{\rA}
  \proofstepwidestar[1]{\SemigroupAlg{A}{\star}}{%
    \FormulaRefAuto{FiniteSemigroupBaseAxiom}[ax]{1}}
  \proofstepwidestar[1,2]{\mathcal R_{A,\star}(a)\neq\varnothing}{%
    \FormulaRefAuto{FiniteSemigroupPowerRepetitionExists}[thm]{1,2}}
  \proofstepwidestar[1,2]{\mathcal R_{A,\star}(a)\subseteq\mathbb N}{%
    \FormulaRefAuto{FiniteSemigroupPowerRepetitionSetDef}[def]{3,2}}
  \proofstepwidestar[1,2]{%
    \exists j\in\mathcal R_{A,\star}(a)\,
      \forall v\in\mathcal R_{A,\star}(a)\,(j\leq v)}{%
    \FormulaRefAuto{NaturalWellOrderingPrinciple}[thm]{5,4}}
  \proofstepwidestar[7]{%
    j\in\mathcal R_{A,\star}(a)\land
      \forall v\in\mathcal R_{A,\star}(a)\,(j\leq v)}{\rA}
  \proofstepwidestar[1,2,7]{j\in\mathbb N}{%
    \FormulaRefAuto{A\subseteq B\dsep x\in A\vdash x\in B}[thm]{%
      5,\rAEa{7}}}
  \proofstepwidestarbreak[1,2,7]{j\neq0}{%
    \FormulaRefAuto{FiniteSemigroupPowerRepetitionSetDef}[def]{%
      3,2,\rAEa{7}}}
  \proofstepwidestarbreak[1,2,7]{%
    \exists i\in\mathbb N\,
      \bigl(i\neq0\land i<j\land a^i=a^j\bigr)}{%
    \FormulaRefAuto{FiniteSemigroupPowerRepetitionSetDef}[def]{%
      3,2,\rAEa{7}}}
  \proofstepwidestar[11]{%
    i\in\mathbb N\land
      \bigl(i\neq0\land i<j\land a^i=a^j\bigr)}{\rA}
  \proofstepwidestarbreak[1,2,7,11]{%
    \exists q\in\mathbb N\,\bigl(i+(q+1)=j\bigr)}{%
    \FormulaRefAuto{PeanoStrictAddOneDistance}[thm]{%
      \rAEa{11},8,\rAEa{\rAEb{\rAEb{11}}}}}
  \proofstepwidestar[13]{q\in\mathbb N\land i+(q+1)=j}{\rA}
  \proofstepwidestarbreak[1,2,7,11,13]{%
    i,j,q\in\mathbb N}{%
    \rAI{\rAEa{11},\rAI{8,\rAEa{13}}}}
  \proofstepwidestar[11]{i\neq0}{\rAEa{\rAEb{11}}}
  \proofstepwidestar[1,2,7]{j\neq0}{\rAIStar{9}}
  \proofstepwidestar[11]{i<j}{%
    \rAEa{\rAEb{\rAEb{11}}}}
  \proofstepwidestarbreak[11]{a^i=a^j}{%
    \rAEb{\rAEb{\rAEb{11}}}}
  \proofstepwidestarbreak[7]{%
    \forall v\in\mathcal R_{A,\star}(a)\,(j\leq v)}{\rAEb{7}}
  \proofstepwidestarbreak[13]{i+(q+1)=j}{\rAEb{13}}
  \proofstepwidestarbreak[1,2,7,11,13]{%
    \operatorname{IP}_{A,\star}(a;i,q+1)}{%
    \FormulaRefAuto{FiniteSemigroupLeastRepetitionProducesIndexPeriod}[thm]{%
      3,2,14,15,16,17,18,19,20}}
  \proofstepwidestar[1,2,7,11,13]{%
    \exists\mu\,\exists\lambda\,
      \operatorname{IP}_{A,\star}(a;\mu,\lambda)}{\rEI{\rEI{21}}}
  \proofstepwidestar[1,2,7,11]{%
    \exists\mu\,\exists\lambda\,
      \operatorname{IP}_{A,\star}(a;\mu,\lambda)}{\rEE{12,13,22}}
  \proofstepwidestar[1,2,7]{%
    \exists\mu\,\exists\lambda\,
      \operatorname{IP}_{A,\star}(a;\mu,\lambda)}{\rEE{10,11,23}}
  \proofstepwidestar[1,2]{%
    \exists\mu\,\exists\lambda\,
      \operatorname{IP}_{A,\star}(a;\mu,\lambda)}{\rEE{6,7,24}}
\end{tabproofwide}


\Needspace{20\baselineskip}
\subsection{Eindeutigkeit von Endpunkt, Index und Periode}
\hypertarget{index-period.proof.FiniteSemigroupIndexPeriodEndpointUnique}{}
\IndexPeriodStatement{FiniteSemigroupIndexPeriodEndpointUnique}
In der Tabelle stehen \(E=\mu+\lambda\) und \(E'=\mu'+\lambda'\).
\begin{tabproofwide}
  \proofstepwidestarbreak[1]{\SemigroupAlg{A}{\star}}{\rA}
  \proofstepwidestarbreak[2]{a\in A}{\rA}
  \proofstepwidestarbreak[3]{\operatorname{IP}_{A,\star}(a;\mu,\lambda)}{\rA}
  \proofstepwidestarbreak[4]{\operatorname{IP}_{A,\star}(a;\mu',\lambda')}{\rA}
  \proofstepwidestarbreak[1--3]{\mu\in\mathbb N\land\mu\neq0\land\lambda\in\mathbb N\land\lambda\neq0}{\FormulaRefAuto{FiniteSemigroupIndexPeriodPositiveData}[ax]{1,2,3}}
  \proofstepwidestarbreak[1,2,4]{\mu'\in\mathbb N\land\mu'\neq0\land\lambda'\in\mathbb N\land\lambda'\neq0}{\FormulaRefAuto{FiniteSemigroupIndexPeriodPositiveData}[ax]{1,2,4}}
  \proofstepwidestarbreak[1--3]{E\in\mathbb N\land\mu<E}{\rAI{\FormulaRefAuto{PeanoAddClosure}[thm]{\rAEa{5},\rAEa{\rAEb{\rAEb{5}}}},\FormulaRefAuto{m\in\mathbb{N}\dsep n\in\mathbb{N}\dsep k\in\mathbb{N}\dsep m+k=n\dsep k\neq0\vdash m<n}[thm]{\rAEa{5},\FormulaRefAuto{PeanoAddClosure}[thm]{\rAEa{5},\rAEa{\rAEb{\rAEb{5}}}},\rAEa{\rAEb{\rAEb{5}}},\rII,\rAEb{\rAEb{\rAEb{5}}}}}}
  \proofstepwidestarbreak[1,2,4]{E'\in\mathbb N\land\mu'<E'}{\rAI{\FormulaRefAuto{PeanoAddClosure}[thm]{\rAEa{6},\rAEa{\rAEb{\rAEb{6}}}},\FormulaRefAuto{m\in\mathbb{N}\dsep n\in\mathbb{N}\dsep k\in\mathbb{N}\dsep m+k=n\dsep k\neq0\vdash m<n}[thm]{\rAEa{6},\FormulaRefAuto{PeanoAddClosure}[thm]{\rAEa{6},\rAEa{\rAEb{\rAEb{6}}}},\rAEa{\rAEb{\rAEb{6}}},\rII,\rAEb{\rAEb{\rAEb{6}}}}}}
  \proofstepwidestarbreak[1--3]{a^E=a^\mu}{\FormulaRefAuto{FiniteSemigroupIndexPeriodEquality}[ax]{1,2,3}}
  \proofstepwidestarbreak[1,2,4]{a^{E'}=a^{\mu'}}{\FormulaRefAuto{FiniteSemigroupIndexPeriodEquality}[ax]{1,2,4}}
  \proofstepwidestarbreak[11]{E<E'}{\rA}
  \proofstepwidestarbreak[1--4,11]{a^\mu\neq a^E}{\FormulaRefAuto{FiniteSemigroupIndexPeriodInitialDistinctness}[ax]{1,2,4,\rAI{\rAEa{5},\rAEa{7}},\rAEa{\rAEb{5}},\rAEb{7},11}}
  \proofstepwidestarbreak[1--4,11]{\bot}{\rBI{12,\FormulaRefAuto{a=b\vdash b=a}[thm]{9}}}
  \proofstepwidestarbreak[1--4]{\neg(E<E')}{\rCI{11,13}}
  \proofstepwidestarbreak[15]{E'<E}{\rA}
  \proofstepwidestarbreak[1--4,15]{a^{\mu'}\neq a^{E'}}{\FormulaRefAuto{FiniteSemigroupIndexPeriodInitialDistinctness}[ax]{1,2,3,\rAI{\rAEa{6},\rAEa{8}},\rAEa{\rAEb{6}},\rAEb{8},15}}
  \proofstepwidestarbreak[1--4,15]{\bot}{\rBI{16,\FormulaRefAuto{a=b\vdash b=a}[thm]{10}}}
  \proofstepwidestarbreak[1--4]{\neg(E'<E)}{\rCI{15,17}}
  \proofstepwidestarbreak[1--4]{E<E'\lor E'<E\lor E=E'}{\FormulaRefAuto{PeanoOrderTrichotomy}[thm]{\rAEa{7},\rAEa{8}}}
  \proofstepwidestarbreak[1--4]{E=E'}{\FormulaRefAuto{P\lor Q,\neg P\vdash Q}[thm]{\FormulaRefAuto{P\lor Q,\neg P\vdash Q}[thm]{19,14},18}}
\end{tabproofwide}


\Needspace{16\baselineskip}
\hypertarget{index-period.proof.FiniteSemigroupIndexPeriodIndexUnique}{}
\IndexPeriodStatement{FiniteSemigroupIndexPeriodIndexUnique}
Wir schreiben \(E=\mu+\lambda=\mu'+\lambda'\).
\begin{tabproofwide}
  \proofstepwidestarbreak[1]{\SemigroupAlg{A}{\star}}{\rA}
  \proofstepwidestarbreak[2]{a\in A}{\rA}
  \proofstepwidestarbreak[3]{\operatorname{IP}_{A,\star}(a;\mu,\lambda)}{\rA}
  \proofstepwidestarbreak[4]{\operatorname{IP}_{A,\star}(a;\mu',\lambda')}{\rA}
  \proofstepwidestarbreak[5]{\mu+\lambda=\mu'+\lambda'}{\rA}
  \proofstepwidestarbreak[1,2,3]{\mu\in\mathbb N\land\mu\neq0\land\lambda\in\mathbb N\land\lambda\neq0}{\FormulaRefAuto{FiniteSemigroupIndexPeriodPositiveData}[ax]{1,2,3}}
  \proofstepwidestarbreak[1,2,4]{\mu'\in\mathbb N\land\mu'\neq0\land\lambda'\in\mathbb N\land\lambda'\neq0}{\FormulaRefAuto{FiniteSemigroupIndexPeriodPositiveData}[ax]{1,2,4}}
  \proofstepwidestarbreak[1,2,3]{\mu<E}{\FormulaRefAuto{m\in\mathbb{N}\dsep n\in\mathbb{N}\dsep k\in\mathbb{N}\dsep m+k=n\dsep k\neq0\vdash m<n}[thm]{\rAEa{6},\FormulaRefAuto{PeanoAddClosure}[thm]{\rAEa{6},\rAEa{\rAEb{\rAEb{6}}}},\rAEa{\rAEb{\rAEb{6}}},\rII,\rAEb{\rAEb{\rAEb{6}}}}}
  \proofstepwidestarbreak[1,2,4,5]{\mu'<E}{\rIE{\FormulaRefAuto{a=b\vdash b=a}[thm]{5},\FormulaRefAuto{m\in\mathbb{N}\dsep n\in\mathbb{N}\dsep k\in\mathbb{N}\dsep m+k=n\dsep k\neq0\vdash m<n}[thm]{\rAEa{7},\FormulaRefAuto{PeanoAddClosure}[thm]{\rAEa{7},\rAEa{\rAEb{\rAEb{7}}}},\rAEa{\rAEb{\rAEb{7}}},\rII,\rAEb{\rAEb{\rAEb{7}}}}}}
  \proofstepwidestarbreak[1,2,3,4,5]{a^\mu=a^{\mu'}}{\rChain{\FormulaRefAuto{a=b\vdash b=a}[thm]{\FormulaRefAuto{FiniteSemigroupIndexPeriodEquality}[ax]{1,2,3}},\rIE{5,\rII},\FormulaRefAuto{FiniteSemigroupIndexPeriodEquality}[ax]{1,2,4}}}
  \proofstepwidestarbreak[11]{\mu<\mu'}{\rA}
  \proofstepwidestarbreak[1,2,3,4,5,11]{a^\mu\neq a^{\mu'}}{\FormulaRefAuto{FiniteSemigroupIndexPeriodInitialDistinctness}[ax]{1,2,3,\rAI{\rAEa{6},\rAEa{7}},\rAEa{\rAEb{6}},11,9}}
  \proofstepwidestarbreak[1,2,3,4,5]{\neg(\mu<\mu')}{\rCI{11,\rBI{12,10}}}
  \proofstepwidestarbreak[14]{\mu'<\mu}{\rA}
  \proofstepwidestarbreak[1,2,3,4,14]{a^{\mu'}\neq a^\mu}{\FormulaRefAuto{FiniteSemigroupIndexPeriodInitialDistinctness}[ax]{1,2,3,\rAI{\rAEa{7},\rAEa{6}},\rAEa{\rAEb{7}},14,8}}
  \proofstepwidestarbreak[1,2,3,4,5]{\neg(\mu'<\mu)}{\rCI{14,\rBI{15,\FormulaRefAuto{a=b\vdash b=a}[thm]{10}}}}
  \proofstepwidestarbreak[1,2,3,4]{\mu<\mu'\lor\mu'<\mu\lor\mu=\mu'}{\FormulaRefAuto{PeanoOrderTrichotomy}[thm]{\rAEa{6},\rAEa{7}}}
  \proofstepwidestarbreak[1,2,3,4,5]{\mu=\mu'}{\FormulaRefAuto{P\lor Q,\neg P\vdash Q}[thm]{\FormulaRefAuto{P\lor Q,\neg P\vdash Q}[thm]{17,13},16}}
\end{tabproofwide}


\hypertarget{index-period.proof.FiniteSemigroupIndexPeriodUniqueness}{}
\IndexPeriodStatement{FiniteSemigroupIndexPeriodUniqueness}
\begin{tabproofwide}
  \proofstepwidestarbreak[1]{\SemigroupAlg{A}{\star}}{\rA}
  \proofstepwidestarbreak[2]{a\in A}{\rA}
  \proofstepwidestarbreak[3]{\operatorname{IP}_{A,\star}(a;\mu,\lambda)}{\rA}
  \proofstepwidestarbreak[4]{\operatorname{IP}_{A,\star}(a;\mu',\lambda')}{\rA}
  \proofstepwidestarbreak[1--4]{\mu+\lambda=\mu'+\lambda'}{\FormulaRefAuto{FiniteSemigroupIndexPeriodEndpointUnique}[thm]{1,2,3,4}}
  \proofstepwidestarbreak[1--4]{\mu=\mu'}{\FormulaRefAuto{FiniteSemigroupIndexPeriodIndexUnique}[thm]{1,2,3,4,5}}
  \proofstepwidestarbreak[1--3]{\mu\in\mathbb N\land\lambda\in\mathbb N}{\rAI{\rAEa{\FormulaRefAuto{FiniteSemigroupIndexPeriodPositiveData}[ax]{1,2,3}},\rAEa{\rAEb{\rAEb{\FormulaRefAuto{FiniteSemigroupIndexPeriodPositiveData}[ax]{1,2,3}}}}}}
  \proofstepwidestarbreak[1,2,4]{\lambda'\in\mathbb N}{\rAEa{\rAEb{\rAEb{\FormulaRefAuto{FiniteSemigroupIndexPeriodPositiveData}[ax]{1,2,4}}}}}
  \proofstepwidestarbreak[1--4]{\mu+\lambda=\mu+\lambda'}{\rIE{\FormulaRefAuto{a=b\vdash b=a}[thm]{6},5}}
  \proofstepwidestarbreak[1--4]{\lambda=\lambda'}{\FormulaRefAuto{PeanoAddLeftCancellation}[thm]{\rAEa{7},\rAEb{7},8,9}}
  \proofstepwidestarbreak[1--4]{\mu=\mu'\land\lambda=\lambda'}{\rAI{6,10}}
\end{tabproofwide}


\Needspace{16\baselineskip}
\hypertarget{index-period.proof.FiniteSemigroupIndexPeriod}{}
\IndexPeriodStatement{FiniteSemigroupIndexPeriod}
Für diese Tabelle bedeutet
\(P(z)\equiv\exists\mu\,\exists\lambda\,
 (z=(\mu,\lambda)\land\operatorname{IP}_{A,\star}(a;\mu,\lambda))\).
\begin{tabproofwide}
  \proofstepwidestarbreak[1]{\FiniteSemigroupAlg{A}{\star}}{\rA}
  \proofstepwidestarbreak[2]{a\in A}{\rA}
  \proofstepwidestarbreak[1]{\SemigroupAlg{A}{\star}}{\FormulaRefAuto{FiniteSemigroupBaseAxiom}[ax]{1}}
  \proofstepwidestarbreak[1,2]{\exists\mu\,\exists\lambda\,\operatorname{IP}_{A,\star}(a;\mu,\lambda)}{\FormulaRefAuto{FiniteSemigroupIndexPeriodExistence}[thm]{1,2}}
  \proofstepwidestarbreak[5]{\exists\lambda\,\operatorname{IP}_{A,\star}(a;\mu,\lambda)}{\rA}
  \proofstepwidestarbreak[6]{\operatorname{IP}_{A,\star}(a;\mu,\lambda)}{\rA}
  \proofstepwidestarbreak[6]{P((\mu,\lambda))}{\rEI{\rEI{\rAI{\rII,6}}}}
  \proofstepwidestarbreak[6]{\exists z\,P(z)}{\rEI{7}}
  \proofstepwidestarbreak[5]{\exists z\,P(z)}{\rEE{5,6,8}}
  \proofstepwidestarbreak[1,2]{\exists z\,P(z)}{\rEE{4,5,9}}
  \proofstepwidestarbreak[11]{P(z)}{\rA}
  \proofstepwidestarbreak[12]{P(z')}{\rA}
  \proofstepwidestarbreak[13]{\exists\lambda\,(z=(\mu,\lambda)\land\operatorname{IP}_{A,\star}(a;\mu,\lambda))}{\rA}
  \proofstepwidestarbreak[14]{z=(\mu,\lambda)\land\operatorname{IP}_{A,\star}(a;\mu,\lambda)}{\rA}
  \proofstepwidestarbreak[15]{\exists\lambda'\,(z'=(\mu',\lambda')\land\operatorname{IP}_{A,\star}(a;\mu',\lambda'))}{\rA}
  \proofstepwidestarbreak[16]{z'=(\mu',\lambda')\land\operatorname{IP}_{A,\star}(a;\mu',\lambda')}{\rA}
  \proofstepwidestarbreak[1,2,14,16]{\mu=\mu'\land\lambda=\lambda'}{\FormulaRefAuto{FiniteSemigroupIndexPeriodUniqueness}[thm]{3,2,\rAEb{14},\rAEb{16}}}
  \proofstepwidestarbreak[1,2,14,16]{z=z'}{\rChain{\rAEa{14},\rIE{\rAEa{17},\rAEb{17},\rII},\FormulaRefAuto{a=b\vdash b=a}[thm]{\rAEa{16}}}}
  \proofstepwidestarbreak[1,2,12,14]{z=z'}{\rEE{12,15,\rEE{15,16,18}}}
  \proofstepwidestarbreak[1,2,11,12]{z=z'}{\rEE{11,13,\rEE{13,14,19}}}
  \proofstepwidestarbreak[1,2]{\exists!z\,P(z)}{\UEI{10,11,12,20}}
\end{tabproofwide}

Die eindeutige Existenz des geordneten Paars rechtfertigt die Bezeichnung
seiner beiden Komponenten als Index und Periode. Die Positivität beider
Komponenten ist bereits eine der bewiesenen Eigenschaften jedes solchen
Paares; es wird keine weitere Auswahlfunktion vorausgesetzt.

\Needspace{24\baselineskip}
\section{Periodische Fortsetzung und die monogene Unterhalbgruppe}
\hypertarget{index-period.proof.FiniteSemigroupPowerEqualityPropagation}{}
\IndexPeriodStatement{FiniteSemigroupPowerEqualityPropagation}
\begin{tabproofwide}
  \proofstepwidestarbreak[1]{\SemigroupAlg{A}{\star}}{\rA}
  \proofstepwidestarbreak[2]{a\in A}{\rA}
  \proofstepwidestarbreak[3]{\mu,\lambda\in\mathbb N}{\rA}
  \proofstepwidestarbreak[4]{\mu\neq0}{\rA}
  \proofstepwidestarbreak[5]{\lambda\neq0}{\rA}
  \proofstepwidestarbreak[6]{a^{\mu+\lambda}=a^\mu}{\rA}
  \proofstepwidestarbreak[7]{t\in\mathbb N}{\rA}
  \proofstepwidestarbreak[3]{\mu+\lambda\in\mathbb N}{\FormulaRefAuto{PeanoAddClosure}[thm]{\rAEa{3},\rAEb{3}}}
  \proofstepwidestarbreak[3,4]{\mu+\lambda\neq0}{\FormulaRefAuto{PeanoAddNonzeroLeft}[thm]{\rAEa{3},\rAEb{3},4}}
  \proofstepwidestarbreak[10]{t=0}{\rA}
  \proofstepwidestarbreak[3,6,7,10]{a^{(\mu+\lambda)+t}=a^{\mu+t}}{\rIE{10,\FormulaRefAuto{PeanoAddRightZero}[thm]{8},\FormulaRefAuto{PeanoAddRightZero}[thm]{\rAEa{3}},6}}
  \proofstepwidestarbreak[12]{t\neq0}{\rA}
  \proofstepwidestarbreak[1--4,7,12]{a^{(\mu+\lambda)+t}=a^{\mu+\lambda}\star a^t}{\FormulaRefAuto{SemigroupPowerAdditionLaw}[thm]{1,2,8,7,9,12}}
  \proofstepwidestarbreak[1--4,6,7,12]{a^{(\mu+\lambda)+t}=a^\mu\star a^t}{\rIE{6,13}}
  \proofstepwidestarbreak[1--4,6,7,12]{a^{(\mu+\lambda)+t}=a^{\mu+t}}{\rChain{14,\FormulaRefAuto{a=b\vdash b=a}[thm]{\FormulaRefAuto{SemigroupPowerAdditionLaw}[thm]{1,2,\rAEa{3},7,4,12}}}}
  \proofstepwidestarbreak[]{t=0\lor t\neq0}{\FormulaRefAuto{P\lor\neg P}[thm]}
  \proofstepwidestarbreak[1--7]{a^{(\mu+\lambda)+t}=a^{\mu+t}}{\rOE{16,10,11,12,15}}
\end{tabproofwide}


\hypertarget{index-period.proof.FiniteSemigroupSinglePeriodShift}{}
\IndexPeriodStatement{FiniteSemigroupSinglePeriodShift}
\begin{tabproofwide}
  \proofstepwidestarbreak[1]{\SemigroupAlg{A}{\star}}{\rA}
  \proofstepwidestarbreak[2]{a\in A}{\rA}
  \proofstepwidestarbreak[3]{\mu,\lambda,r\in\mathbb N}{\rA}
  \proofstepwidestarbreak[4]{\mu\neq0}{\rA}
  \proofstepwidestarbreak[5]{\lambda\neq0}{\rA}
  \proofstepwidestarbreak[6]{a^{\mu+\lambda}=a^\mu}{\rA}
  \proofstepwidestarbreak[7]{\mu\leq r}{\rA}
  \proofstepwidestarbreak[3,7]{\exists t\in\mathbb N\,(\mu+t=r)}{\FormulaRefAuto{NaturalOrderAdditiveCharacterization}[thm]{\rAEa{3},\rAEb{\rAEb{3}},7}}
  \proofstepwidestarbreak[9]{t\in\mathbb N\land\mu+t=r}{\rA}
  \proofstepwidestarbreak[3,9]{(\mu+t)+\lambda=(\mu+\lambda)+t}{\rChain{\FormulaRefAuto{PeanoAddAssociative}[thm]{\rAEa{3},\rAEa{9},\rAEa{\rAEb{3}}},\rIE{\FormulaRefAuto{PeanoAddCommutative}[thm]{\rAEa{9},\rAEa{\rAEb{3}}}},\FormulaRefAuto{a=b\vdash b=a}[thm]{\FormulaRefAuto{PeanoAddAssociative}[thm]{\rAEa{3},\rAEa{\rAEb{3}},\rAEa{9}}}}}
  \proofstepwidestarbreak[1--6,9]{a^{(\mu+\lambda)+t}=a^{\mu+t}}{\FormulaRefAuto{FiniteSemigroupPowerEqualityPropagation}[thm]{1,2,\rAI{\rAEa{3},\rAEa{\rAEb{3}}},4,5,6,\rAEa{9}}}
  \proofstepwidestarbreak[1--6,9]{a^{r+\lambda}=a^r}{\rIE{\rAEb{9},10,11}}
  \proofstepwidestarbreak[1--7]{a^{r+\lambda}=a^r}{\rEE{8,9,12}}
\end{tabproofwide}


\hypertarget{index-period.proof.FiniteSemigroupPowerEventuallyPeriodic}{}
\IndexPeriodStatement{FiniteSemigroupPowerEventuallyPeriodic}
\begin{tabproofwide}
  \proofstepwidestarbreak[1]{\SemigroupAlg{A}{\star}}{\rA}
  \proofstepwidestarbreak[2]{a\in A}{\rA}
  \proofstepwidestarbreak[3]{\operatorname{IP}_{A,\star}(a;\mu,\lambda)}{\rA}
  \proofstepwidestarbreak[1--3]{\mu\in\mathbb N\land\mu\neq0\land\lambda\in\mathbb N\land\lambda\neq0}{\FormulaRefAuto{FiniteSemigroupIndexPeriodPositiveData}[ax]{1,2,3}}
  \proofstepwidestarbreak[1--3]{a^{\mu+\lambda}=a^\mu}{\FormulaRefAuto{FiniteSemigroupIndexPeriodEquality}[ax]{1,2,3}}
  \proofstepwidestarbreak[6]{r\in\mathbb N}{\rA}
  \proofstepwidestarbreak[7]{\mu\leq r}{\rA}
  \proofstepwidestarbreak[1--3,6,7]{a^{r+\lambda}=a^r}{\FormulaRefAuto{FiniteSemigroupSinglePeriodShift}[thm]{1,2,\rAI{\rAEa{4},\rAI{\rAEa{\rAEb{\rAEb{4}}},6}},\rAEa{\rAEb{4}},\rAEb{\rAEb{\rAEb{4}}},5,7}}
  \proofstepwidestarbreak[1--3]{\forall r\in\mathbb N\,(\mu\leq r\rightarrow a^{r+\lambda}=a^r)}{\rUI{\rRI{6,\rRI{7,8}}}}
\end{tabproofwide}


\Needspace{20\baselineskip}
\subsection{Reduktion und Anzahl der verschiedenen Potenzen}
\hypertarget{index-period.proof.FiniteSemigroupMonogenicSize}{}
\IndexPeriodStatement{FiniteSemigroupMonogenicSize}
Die Beweistabelle trennt die Reduktion beliebiger Potenzen von der
Zählung der verschiedenen Anfangswerte. Wir schreiben
\[
 \begin{aligned}
 E&=\mu+\lambda,\\
 P(r)&\equiv\exists s\in\mathbb N\,
       (s\neq0\land s<E\land a^r=a^s),\\
 M&=\{r\in\mathbb N\mid r\neq0\land\neg P(r)\}.
 \end{aligned}
\]
\begin{tabproofsplitwide}
\noindent\textbf{Aussage des ersten Beweisteils.}
\[
 \begin{aligned}[t]
 &\SemigroupAlg{A}{\star}\dsep a\in A\dsep
 \operatorname{IP}_{A,\star}(a;\mu,\lambda)\vdash{}\\[-2pt]
 &\forall r\,\bigl(r\in\mathbb N\land r\neq0\rightarrow{}\\[-2pt]
 &\qquad\exists s\in\mathbb N\,(s\neq0\land s<\mu+\lambda\land a^r=a^s)\bigr).
 \end{aligned}
\]
\proofpartwide{Reduktion auf den ersten Potenzabschnitt}
  \proofstepwidestarbreak[1]{\SemigroupAlg{A}{\star}}{\rA}
  \proofstepwidestarbreak[2]{a\in A}{\rA}
  \proofstepwidestarbreak[3]{\operatorname{IP}_{A,\star}(a;\mu,\lambda)}{\rA}
  \proofstepwidestarbreak[1,2,3]{\mu\in\mathbb N\land\mu\neq0\land\lambda\in\mathbb N\land\lambda\neq0}{\FormulaRefAuto{FiniteSemigroupIndexPeriodPositiveData}[ax]{1,2,3}}
  \proofstepwidestarbreak[1,2,3]{E\in\mathbb N}{\FormulaRefAuto{PeanoAddClosure}[thm]{\rAEa{4},\rAEa{\rAEb{\rAEb{4}}}}}
  \proofstepwidestarbreak[1,2,3]{a^E=a^\mu}{\FormulaRefAuto{FiniteSemigroupIndexPeriodEquality}[ax]{1,2,3}}
  \proofstepwidestarbreak[]{M\subseteq\mathbb N}{\FormulaRefAuto{\{x\in A\mid P(x)\}\subseteq A}[thm]}
  \proofstepwidestarbreak[8]{M\neq\varnothing}{\rA}
  \proofstepwidestarbreak[8]{\exists r\in M\,\forall v\in M\,(r\leq v)}{\FormulaRefAuto{NaturalWellOrderingPrinciple}[thm]{7,8}}
  \proofstepwidestarbreak[10]{r\in M\land\forall v\in M\,(r\leq v)}{\rA}
  \proofstepwidestarbreak[10]{r\in\mathbb N\land r\neq0\land\neg P(r)}{\FormulaRefAuto{x\in\{u\in A\mid P(u)\}\eqvdash x\in A\land P(x)}[thm]{\rAEa{10}}}
  \proofstepwidestarbreak[12]{r<E}{\rA}
  \proofstepwidestarbreak[10,12]{P(r)}{\rEI{\rAI{\rAEa{11},\rAI{\rAEa{\rAEb{11}},\rAI{12,\rII}}}}}
  \proofstepwidestarbreak[10]{\neg(r<E)}{\rCI{12,\rBI{\rAEb{\rAEb{11}},13}}}
  \proofstepwidestarbreak[1,2,3,10]{E<r\lor r<E\lor E=r}{\FormulaRefAuto{PeanoOrderTrichotomy}[thm]{5,\rAEa{11}}}
  \proofstepwidestarbreak[16]{E<r}{\rA}
  \proofstepwidestarbreak[1,2,3,10,16]{E\leq r}{\FormulaRefAuto{PeanoStrictImpliesLeq}[thm]{5,\rAEa{11},16}}
  \proofstepwidestarbreak[18]{r<E\lor E=r}{\rA}
  \proofstepwidestarbreak[19]{r<E}{\rA}
  \proofstepwidestarbreak[10,19]{E\leq r}{\FormulaRefAuto{\bot\vdash P}[thm]{\rBI{14,19}}}
  \proofstepwidestarbreak[21]{E=r}{\rA}
  \proofstepwidestarbreak[1,2,3,21]{E\leq r}{\rIE{21,\FormulaRefAuto{PeanoLeqReflexive}[thm]{5}}}
  \proofstepwidestarbreak[1,2,3,10]{E\leq r}{\rOE{15,16,17,18,\rOE{18,19,20,21,22}}}
  \proofstepwidestarbreak[1,2,3,10]{\exists t\in\mathbb N\,(E+t=r)}{\FormulaRefAuto{NaturalOrderAdditiveCharacterization}[thm]{5,\rAEa{11},23}}
  \proofstepwidestarbreak[25]{t\in\mathbb N\land E+t=r}{\rA}
  \proofstepwidestarbreak[1,2,3,25]{\mu+t\in\mathbb N}{\FormulaRefAuto{PeanoAddClosure}[thm]{\rAEa{4},\rAEa{25}}}
  \proofstepwidestarbreak[1,2,3,25]{\mu+t\neq0}{\FormulaRefAuto{PeanoAddNonzeroLeft}[thm]{\rAEa{4},\rAEa{25},\rAEa{\rAEb{4}}}}
  \proofstepwidestarbreak[1,2,3,25]{(\mu+t)+\lambda=r}{\rChain{\FormulaRefAuto{PeanoAddThreeTermSwap}[thm]{\rAEa{4},\rAEa{25},\rAEa{\rAEb{\rAEb{4}}}},\rAEb{25}}}
  \proofstepwidestarbreak[1,2,3,10,25]{\mu+t<r}{\FormulaRefAuto{m\in\mathbb{N}\dsep n\in\mathbb{N}\dsep k\in\mathbb{N}\dsep m+k=n\dsep k\neq0\vdash m<n}[thm]{26,\rAEa{11},\rAEa{\rAEb{\rAEb{4}}},28,\rAEb{\rAEb{\rAEb{4}}}}}
  \proofstepwidestarbreak[1,2,3,25]{a^r=a^{\mu+t}}{\rIE{\rAEb{25},\FormulaRefAuto{FiniteSemigroupPowerEqualityPropagation}[thm]{1,2,\rAI{\rAEa{4},\rAEa{\rAEb{\rAEb{4}}}},\rAEa{\rAEb{4}},\rAEb{\rAEb{\rAEb{4}}},6,\rAEa{25}}}}
  \proofstepwidestarbreak[31]{\neg P(\mu+t)}{\rA}
  \proofstepwidestarbreak[1,2,3,25,31]{\mu+t\in M}{\FormulaRefAuto{x\in\{u\in A\mid P(u)\}\eqvdash x\in A\land P(x)}[thm]{\rAI{26,\rAI{27,31}}}}
  \proofstepwidestarbreak[1,2,3,10,25,31]{r\leq\mu+t}{\rRE{32,\rUE{\rAEb{10}}}}
  \proofstepwidestarbreak[1,2,3,10,25,31]{\bot}{\FormulaRefAuto{PeanoOrderNoCycle}[thm]{26,\rAEa{11},29,33}}
  \proofstepwidestarbreak[1,2,3,10,25]{P(\mu+t)}{\rCE{31,34}}
  \proofstepwidestarbreak[1,2,3,10,25]{P(r)}{\rIE{\FormulaRefAuto{a=b\vdash b=a}[thm]{30},35}}
  \proofstepwidestarbreak[1,2,3,10,25]{\bot}{\rBI{\rAEb{\rAEb{11}},36}}
  \proofstepwidestarbreak[1,2,3,10]{\bot}{\rEE{24,25,37}}
  \proofstepwidestarbreak[1,2,3,8]{\bot}{\rEE{9,10,38}}
  \proofstepwidestarbreak[1,2,3]{\neg(M\neq\varnothing)}{\rCI{8,39}}
  \proofstepwidestarbreak[41]{r\in\mathbb N\land r\neq0}{\rA}
  \proofstepwidestarbreak[42]{\neg P(r)}{\rA}
  \proofstepwidestarbreak[41,42]{r\in M}{\FormulaRefAuto{x\in\{u\in A\mid P(u)\}\eqvdash x\in A\land P(x)}[thm]{\rAI{\rAEa{41},\rAI{\rAEb{41},42}}}}
  \proofstepwidestarbreak[41,42]{M\neq\varnothing}{\FormulaRefAuto{x\in A\vdash A\neq\varnothing}[thm]{43}}
  \proofstepwidestarbreak[1,2,3,41]{P(r)}{\rCE{42,\rBI{40,44}}}
  \proofstepwidestarbreak[1,2,3]{\forall r\,(r\in\mathbb N\land r\neq0\rightarrow P(r))}{\rUI{\rRI{41,45}}}
\closeproofpartwide

Für die Zählung verwenden wir die Potenzfolge \(\sigma\) aus Band~28 und
ihre Einschränkung \(f\) auf einen endlichen Anfangsabschnitt:
\[
 \begin{gathered}
 n=E-1,\qquad D=\NatSeg n,\qquad
 \sigma=\operatorname{pow}_{A,\star}(a),\qquad f=\sigma\restriction_D,\\
 B=\{y\in A\mid\exists r\in\mathbb N\,
       (r\neq0\land r<E\land y=a^r)\}.
 \end{gathered}
\]
Die eben bewiesene Aussage wird im zweiten Beweisteil als \(R\) abgekürzt.
\noindent\textbf{Aussage des zweiten Beweisteils.}
\[
 \begin{aligned}[t]
 &\SemigroupAlg{A}{\star}\dsep a\in A\dsep
 \operatorname{IP}_{A,\star}(a;\mu,\lambda)\vdash{}\\[-2pt]
 &\langle a\rangle_{A,\star}=B\land\Finite{\langle a\rangle_{A,\star}}\\[-2pt]
 &\qquad{}\land\FiniteCard(\langle a\rangle_{A,\star})=\mu+\lambda-1.
 \end{aligned}
\]
\proofpartwide{Endlichkeit und Kardinalität}
\setcounter{proofstepnr}{46}
  \proofstepwidestarbreak[1,2,3]{R}{\rAIStar{46}}
  \proofstepwidestarbreak[1,2,3]{\mu\in\mathbb N\land\mu\neq0\land\lambda\in\mathbb N\land\lambda\neq0}{\FormulaRefAuto{FiniteSemigroupIndexPeriodPositiveData}[ax]{1,2,3}}
  \proofstepwidestarbreak[1,2,3]{E\in\mathbb N\land E\neq0}{\rAI{\FormulaRefAuto{PeanoAddClosure}[thm]{\rAEa{48},\rAEa{\rAEb{\rAEb{48}}}},\FormulaRefAuto{PeanoAddNonzeroLeft}[thm]{\rAEa{48},\rAEa{\rAEb{\rAEb{48}}},\rAEa{\rAEb{48}}}}}
  \proofstepwidestarbreak[1,2,3]{n\in\mathbb N}{\FormulaRefAuto{PeanoMinusOneNatural}[thm]{\rAEa{49},\rAEb{49}}}
  \proofstepwidestarbreak[1,2,3]{E=n+1}{\FormulaRefAuto{PeanoMinusOneRecoverValue}[thm]{\rAEa{49},\rAEb{49}}}
  \proofstepwidestarbreak[1,2,3]{D\subseteq\mathbb N\land\Finite D}{\rAI{\FormulaRefAuto{PeanoNatSegSubsetN}[thm]{50},\FormulaRefAuto{PeanoNatSegFinite}[thm]{50}}}
  \proofstepwidestarbreak[1,2]{\sigma:\mathbb N\to A}{\FormulaRefAuto{SemigroupPowerSequenceFunction}[thm]{1,2}}
  \proofstepwidestarbreak[1,2,3]{f:D\to A}{\FormulaRefAuto{C\subseteq A\vdash F\restriction_C\colon C\to B}[thm]{\rAEa{52},53}}
  \proofstepwidestarbreak[55]{k\in D}{\rA}
  \proofstepwidestarbreak[1,2,3,55]{k\in\mathbb N\land k<n}{\rAI{\FormulaRefAuto{PeanoNatSegElementNatural}[thm]{50,55},\FormulaRefAuto{PeanoNatSegMembership}[thm]{50,\FormulaRefAuto{PeanoNatSegElementNatural}[thm]{50,55},55}}}
  \proofstepwidestarbreak[1,2,3,55]{k+1\in\mathbb N\land k+1\neq0}{\rAI{\FormulaRefAuto{PeanoAddOneClosure}[thm]{\rAEa{56}},\FormulaRefAuto{PeanoAddOneNonzero}[thm]{\rAEa{56}}}}
  \proofstepwidestarbreak[1,2,3,55]{k+1<E}{\rIE{\FormulaRefAuto{a=b\vdash b=a}[thm]{\FormulaRefAuto{PeanoAddCommutative}[thm]{\rAEa{56},\FormulaRefAuto{PeanoOneInN}[thm]}},\FormulaRefAuto{a=b\vdash b=a}[thm]{\FormulaRefAuto{PeanoAddCommutative}[thm]{50,\FormulaRefAuto{PeanoOneInN}[thm]}},\FormulaRefAuto{a=b\vdash b=a}[thm]{51},\FormulaRefAuto{PeanoAddLeftStrictMonotone}[thm]{\FormulaRefAuto{PeanoOneInN}[thm],\rAEa{56},50,\rAEb{56}}}}
  \proofstepwidestarbreak[1,2,3,55]{f(k)=a^{k+1}}{\rChain{\FormulaRefAuto{RestrictionDef}[def]{\rAEa{52},53,55},\FormulaRefAuto{a=b\vdash b=a}[thm]{\FormulaRefAuto{SemigroupPowerSuccessorIndex}[thm]{1,2,\rAEa{56}}}}}
  \proofstepwidestarbreak[1,2,3,55]{f(k)\in B}{\FormulaRefAuto{x\in\{u\in A\mid P(u)\}\eqvdash x\in A\land P(x)}[thm]{\rAI{\FormulaRefAuto{F\colon A\to B\dsep x\in A\vdash F(x)\in B}[thm]{54,55},\rEI{\rAI{\rAEa{57},\rAI{\rAEb{57},\rAI{58,59}}}}}}}
  \proofstepwidestarbreak[1,2,3]{\forall k\in D\,(f(k)\in B)}{\rUI{\rRI{55,60}}}
  \proofstepwidestarbreak[62]{r\in\mathbb N\land r\neq0\land r<E}{\rA}
  \proofstepwidestarbreak[62]{r-1\in\mathbb N}{\FormulaRefAuto{PeanoMinusOneNatural}[thm]{\rAEa{62},\rAEa{\rAEb{62}}}}
  \proofstepwidestarbreak[1,2,3,62]{r-1<n}{\FormulaRefAuto{PeanoMinusOneStrictMonotonePositive}[thm]{\rAI{\rAEa{62},\rAEa{49}},\rAEa{\rAEb{62}},\rAEb{\rAEb{62}}}}
  \proofstepwidestarbreak[1,2,3,62]{r-1\in D}{\FormulaRefAuto{PeanoNatSegMembership}[thm]{50,63,64}}
  \proofstepwidestarbreak[62]{r=(r-1)+1}{\FormulaRefAuto{PeanoMinusOneRecoverValue}[thm]{\rAEa{62},\rAEa{\rAEb{62}}}}
  \proofstepwidestarbreak[1,2,3,62]{f(r-1)=a^r}{\rIE{\FormulaRefAuto{a=b\vdash b=a}[thm]{66},\rRE{65,\rUE{\rUI{\rRI{55,59}}}}}}
  \proofstepwidestarbreak[1,2,3]{\forall r\,(r\in\mathbb N\land r\neq0\land r<E\rightarrow\exists k\in D\,(f(k)=a^r))}{\rUI{\rRI{62,\rEI{\rAI{65,67}}}}}
  \proofstepwidestarbreak[69]{y\in f[D]}{\rA}
  \proofstepwidestarbreak[1,2,3,69]{\exists k\in D\,(y=f(k))}{\FormulaRefAuto{C\subseteq A\dsep F\colon A\to B\dsep y\in F[C]\vdash\exists x\in C\,y=F(x)}[thm]{\FormulaRefAuto{A\subseteq A}[thm],54,69}}
  \proofstepwidestarbreak[71]{k\in D\land y=f(k)}{\rA}
  \proofstepwidestarbreak[1,2,3,71]{f(k)\in B}{\rRE{\rAEa{71},\rUE{61}}}
  \proofstepwidestarbreak[1,2,3,71]{y\in B}{\rIE{\FormulaRefAuto{a=b\vdash b=a}[thm]{\rAEb{71}},72}}
  \proofstepwidestarbreak[1,2,3,69]{y\in B}{\rEE{70,71,73}}
  \proofstepwidestarbreak[1,2,3]{f[D]\subseteq B}{\FormulaRefAuto{A\subseteq B\coloneq\forall x\,(x\in A\to x\in B)}[def]{\rUI{\rRI{69,74}}}}
  \proofstepwidestarbreak[76]{y\in B}{\rA}
  \proofstepwidestarbreak[76]{y\in A\land\exists r\in\mathbb N\,(r\neq0\land r<E\land y=a^r)}{\FormulaRefAuto{x\in\{u\in A\mid P(u)\}\eqvdash x\in A\land P(x)}[thm]{76}}
  \proofstepwidestarbreak[78]{r\in\mathbb N\land r\neq0\land r<E\land y=a^r}{\rA}
  \proofstepwidestarbreak[1,2,3,78]{\exists k\in D\,(f(k)=a^r)}{\rRE{\rAI{\rAEa{78},\rAI{\rAEa{\rAEb{78}},\rAEa{\rAEb{\rAEb{78}}}}},\rUE{68}}}
  \proofstepwidestarbreak[80]{k\in D\land f(k)=a^r}{\rA}
  \proofstepwidestarbreak[78,80]{y=f(k)}{\rChain{\rAEb{\rAEb{\rAEb{78}}},\FormulaRefAuto{a=b\vdash b=a}[thm]{\rAEb{80}}}}
  \proofstepwidestarbreak[1,2,3,78,80]{y\in f[D]}{\FormulaRefAuto{C\subseteq A\dsep F\colon A\to B\vdash F[C]\coloneqq\iota D\bigl(\forall y\,(y\in D\leftrightarrow\exists x\in C\,y=F(x))\bigr)}[def]{\FormulaRefAuto{A\subseteq A}[thm],54,\rEI{\rAI{\rAEa{80},81}}}}
  \proofstepwidestarbreak[1,2,3,78]{y\in f[D]}{\rEE{79,80,82}}
  \proofstepwidestarbreak[1,2,3,76]{y\in f[D]}{\rEE{\rAEb{77},78,83}}
  \proofstepwidestarbreak[1,2,3]{B\subseteq f[D]}{\FormulaRefAuto{A\subseteq B\coloneq\forall x\,(x\in A\to x\in B)}[def]{\rUI{\rRI{76,84}}}}
  \proofstepwidestarbreak[1,2,3]{f[D]=B}{\FormulaRefAuto{A\subseteq B\land B\subseteq A\eqvdash A=B}[thm]{\rAI{75,85}}}
  \proofstepwidestarbreak[87]{k\in D}{\rA}
  \proofstepwidestarbreak[88]{l\in D}{\rA}
  \proofstepwidestarbreak[89]{f(k)=f(l)}{\rA}
  \proofstepwidestarbreak[1,2,3,87,88]{\begin{aligned}[t]&(k+1\in\mathbb N\land k+1\neq0)\land{}\\[-2pt]&\bigl((l+1\in\mathbb N\land l+1\neq0)\land{}\\[-2pt]&\quad(k+1<E\land l+1<E)\bigr)\end{aligned}}{\rAI{\rRE{87,\rUE{\rUI{\rRI{55,57}}}},\rAI{\rRE{88,\rUE{\rUI{\rRI{55,57}}}},\rAI{\rRE{87,\rUE{\rUI{\rRI{55,58}}}},\rRE{88,\rUE{\rUI{\rRI{55,58}}}}}}}}
  \proofstepwidestarbreak[1,2,3,87,88,89]{a^{k+1}=a^{l+1}}{\rIE{\rRE{87,\rUE{\rUI{\rRI{55,59}}}},\rRE{88,\rUE{\rUI{\rRI{55,59}}}},89}}
  \proofstepwidestarbreak[92]{k+1<l+1}{\rA}
  \proofstepwidestarbreak[1,2,3,87,88,89,92]{\bot}{\rBI{\FormulaRefAuto{FiniteSemigroupIndexPeriodInitialDistinctness}[ax]{1,2,3,\rAI{\rAEa{\rAEa{90}},\rAEa{\rAEa{\rAEb{90}}}},\rAEb{\rAEa{90}},92,\rAEb{\rAEb{\rAEb{90}}}},91}}
  \proofstepwidestarbreak[1,2,3,87,88,89]{\neg(k+1<l+1)}{\rCI{92,93}}
  \proofstepwidestarbreak[95]{l+1<k+1}{\rA}
  \proofstepwidestarbreak[1,2,3,87,88,89,95]{\bot}{\rBI{\FormulaRefAuto{FiniteSemigroupIndexPeriodInitialDistinctness}[ax]{1,2,3,\rAI{\rAEa{\rAEa{\rAEb{90}}},\rAEa{\rAEa{90}}},\rAEb{\rAEa{\rAEb{90}}},95,\rAEa{\rAEb{\rAEb{90}}}},\FormulaRefAuto{a=b\vdash b=a}[thm]{91}}}
  \proofstepwidestarbreak[1,2,3,87,88,89]{\neg(l+1<k+1)}{\rCI{95,96}}
  \proofstepwidestarbreak[1,2,3,87,88,89]{k+1=l+1}{\FormulaRefAuto{P\lor Q,\neg P\vdash Q}[thm]{\FormulaRefAuto{P\lor Q,\neg P\vdash Q}[thm]{\FormulaRefAuto{PeanoOrderTrichotomy}[thm]{\rAEa{\rAEa{90}},\rAEa{\rAEa{\rAEb{90}}}},94},97}}
  \proofstepwidestarbreak[1,2,3,87,88,89]{k=l}{\FormulaRefAuto{PeanoAddOneInjective}[thm]{\FormulaRefAuto{PeanoNatSegElementNatural}[thm]{50,87},\FormulaRefAuto{PeanoNatSegElementNatural}[thm]{50,88},98}}
  \proofstepwidestarbreak[1,2,3]{f:D\inj A}{\FormulaRefAuto{Injektive Funktion}[def]{54,\rUI{\rRI{87,\rUI{\rRI{88,\rRI{89,99}}}}}}}
  \proofstepwidestarbreak[1,2,3]{f:D\bij f[D]}{\FormulaRefAuto{F\colon A\inj B\vdash F\colon A\bij F[A]}[thm]{100}}
  \proofstepwidestarbreak[1,2,3]{D\EqCard B}{\FormulaRefAuto{Gleichmächtigkeit}[def]{\rEI{\rIE{86,101}}}}
  \proofstepwidestarbreak[1,2,3]{\Finite B}{\FormulaRefAuto{FiniteEqCardTransport}[thm]{\rAEb{52},102}}
  \proofstepwidestarbreak[1,2,3]{\FiniteCard(B)=n}{\FormulaRefAuto{FiniteCardWitnessUnique}[thm]{\FormulaRefAuto{FiniteCardNatural}[thm]{103},50,\FormulaRefAuto{FiniteCardCounts}[thm]{103},102}}
  \proofstepwidestarbreak[1,2]{\langle a\rangle_{A,\star}=P_{A,\star}(a)}{\FormulaRefAuto{SemigroupMonogenicPowerDescription}[thm]{1,2}}
  \proofstepwidestarbreak[106]{y\in\langle a\rangle_{A,\star}}{\rA}
  \proofstepwidestarbreak[1,2,106]{y\in P_{A,\star}(a)}{\rIE{105,106}}
  \proofstepwidestarbreak[1,2,106]{y\in A\land\exists r\in\mathbb N\,(r\neq0\land y=a^r)}{\FormulaRefAuto{SemigroupPositivePowerSetMembership}[thm]{1,2,107}}
  \proofstepwidestarbreak[109]{r\in\mathbb N\land r\neq0\land y=a^r}{\rA}
  \proofstepwidestarbreak[1,2,3,109]{P(r)}{\rRE{\rAI{\rAEa{109},\rAEa{\rAEb{109}}},\rUE{47}}}
  \proofstepwidestarbreak[111]{s\in\mathbb N\land s\neq0\land s<E\land a^r=a^s}{\rA}
  \proofstepwidestarbreak[109,111]{y=a^s}{\rChain{\rAEb{\rAEb{109}},\rAEb{\rAEb{\rAEb{111}}}}}
  \proofstepwidestarbreak[1,2,106,109,111]{y\in B}{\FormulaRefAuto{x\in\{u\in A\mid P(u)\}\eqvdash x\in A\land P(x)}[thm]{\rAI{\rAEa{108},\rEI{\rAI{\rAEa{111},\rAI{\rAEa{\rAEb{111}},\rAI{\rAEa{\rAEb{\rAEb{111}}},112}}}}}}}
  \proofstepwidestarbreak[1,2,3,106,109]{y\in B}{\rEE{110,111,113}}
  \proofstepwidestarbreak[1,2,3,106]{y\in B}{\rEE{\rAEb{108},109,114}}
  \proofstepwidestarbreak[1,2,3]{\langle a\rangle_{A,\star}\subseteq B}{\FormulaRefAuto{A\subseteq B\coloneq\forall x\,(x\in A\to x\in B)}[def]{\rUI{\rRI{106,115}}}}
  \proofstepwidestarbreak[117]{y\in B}{\rA}
  \proofstepwidestarbreak[117]{y\in A\land\exists s\in\mathbb N\,(s\neq0\land s<E\land y=a^s)}{\FormulaRefAuto{x\in\{u\in A\mid P(u)\}\eqvdash x\in A\land P(x)}[thm]{117}}
  \proofstepwidestarbreak[119]{s\in\mathbb N\land s\neq0\land s<E\land y=a^s}{\rA}
  \proofstepwidestarbreak[1,2,117,119]{y\in P_{A,\star}(a)}{\FormulaRefAuto{SemigroupPositivePowerSetMembership}[thm]{1,2,\rAI{\rAEa{118},\rEI{\rAI{\rAEa{119},\rAI{\rAEa{\rAEb{119}},\rAEb{\rAEb{\rAEb{119}}}}}}}}}
  \proofstepwidestarbreak[1,2,117,119]{y\in\langle a\rangle_{A,\star}}{\rIE{\FormulaRefAuto{a=b\vdash b=a}[thm]{105},120}}
  \proofstepwidestarbreak[1,2,117]{y\in\langle a\rangle_{A,\star}}{\rEE{\rAEb{118},119,121}}
  \proofstepwidestarbreak[1,2]{B\subseteq\langle a\rangle_{A,\star}}{\FormulaRefAuto{A\subseteq B\coloneq\forall x\,(x\in A\to x\in B)}[def]{\rUI{\rRI{117,122}}}}
  \proofstepwidestarbreak[1,2,3]{\langle a\rangle_{A,\star}=B}{\FormulaRefAuto{A\subseteq B\land B\subseteq A\eqvdash A=B}[thm]{\rAI{116,123}}}
  \proofstepwidestarbreak[1,2,3]{\begin{aligned}[t]&\langle a\rangle_{A,\star}=B\\[-2pt]&{}\land\Finite{\langle a\rangle_{A,\star}}\\[-2pt]&{}\land\FiniteCard(\langle a\rangle_{A,\star})=\mu+\lambda-1\end{aligned}}{\rAI{124,\rIE{\FormulaRefAuto{a=b\vdash b=a}[thm]{124},\rAI{103,104}}}}
\closeproofpartwide
\end{tabproofsplitwide}


\Needspace{24\baselineskip}
\section{Idempotente Potenzen}
\hypertarget{index-period.proof.FiniteSemigroupMultiplePeriodShift}{}
\IndexPeriodStatement{FiniteSemigroupMultiplePeriodShift}
Die Induktion läuft über die Anzahl der Perioden \(q\).
\begin{tabproofwide}
  \proofstepwidestarbreak[1]{\SemigroupAlg{A}{\star}}{\rA}
  \proofstepwidestarbreak[2]{a\in A}{\rA}
  \proofstepwidestarbreak[3]{\mu,\lambda,r,q\in\mathbb N}{\rA}
  \proofstepwidestarbreak[4]{\mu\neq0}{\rA}
  \proofstepwidestarbreak[5]{\lambda\neq0}{\rA}
  \proofstepwidestarbreak[6]{a^{\mu+\lambda}=a^\mu}{\rA}
  \proofstepwidestarbreak[7]{\mu\leq r}{\rA}
  \proofstepwidestarbreak[3]{r+\lambda\cdot0=r}{\rChain{\rIE{\FormulaRefAuto{PeanoMulRightZero}[thm]{\rAEa{\rAEb{3}}},\rII},\FormulaRefAuto{PeanoAddRightZero}[thm]{\rAEa{\rAEb{\rAEb{3}}}}}}
  \proofstepwidestarbreak[3]{a^{r+\lambda\cdot0}=a^r}{\rIE{8,\rII}}
  \proofstepwidestarbreak[10]{k\in\mathbb N}{\rA}
  \proofstepwidestarbreak[11]{a^{r+\lambda\cdot k}=a^r}{\rA}
  \proofstepwidestarbreak[3,10]{\lambda\cdot k\in\mathbb N}{\FormulaRefAuto{PeanoMulClosure}[thm]{\rAEa{\rAEb{3}},10}}
  \proofstepwidestarbreak[3,10]{r+\lambda\cdot k\in\mathbb N}{\FormulaRefAuto{PeanoAddClosure}[thm]{\rAEa{\rAEb{\rAEb{3}}},12}}
  \proofstepwidestarbreak[3,10]{r\leq r+\lambda\cdot k}{\FormulaRefAuto{NaturalOrderAdditiveCharacterization}[thm]{\rAEa{\rAEb{\rAEb{3}}},13,\rEI{\rAI{12,\rII}}}}
  \proofstepwidestarbreak[3,7,10]{\mu\leq r+\lambda\cdot k}{\FormulaRefAuto{PeanoLeqTransitive}[thm]{\rAEa{3},\rAEa{\rAEb{\rAEb{3}}},13,7,14}}
  \proofstepwidestarbreak[1,2,3,4,5,6,7,10]{a^{(r+\lambda\cdot k)+\lambda}=a^{r+\lambda\cdot k}}{\FormulaRefAuto{FiniteSemigroupSinglePeriodShift}[thm]{1,2,\rAI{\rAEa{3},\rAI{\rAEa{\rAEb{3}},13}},4,5,6,15}}
  \proofstepwidestarbreak[3,10]{r+\lambda\cdot(k+1)=(r+\lambda\cdot k)+\lambda}{\rChain{\rIE{\FormulaRefAuto{PeanoMulAddOneRight}[thm]{\rAEa{\rAEb{3}},10},\rII},\FormulaRefAuto{a=b\vdash b=a}[thm]{\FormulaRefAuto{PeanoAddAssociative}[thm]{\rAEa{\rAEb{\rAEb{3}}},12,\rAEa{\rAEb{3}}}}}}
  \proofstepwidestarbreak[1,2,3,4,5,6,7,10,11]{a^{r+\lambda\cdot(k+1)}=a^r}{\rChain{\rIE{17,\rII},16,11}}
  \proofstepwidestarbreak[1,2,3,4,5,6,7]{\forall k\in\mathbb N\,(a^{r+\lambda\cdot k}=a^r\rightarrow a^{r+\lambda\cdot(k+1)}=a^r)}{\rUI{\rRI{10,\rRI{11,18}}}}
  \proofstepwidestarbreak[1,2,3,4,5,6,7]{\forall q\in\mathbb N\,(a^{r+\lambda\cdot q}=a^r)}{\FormulaRefAuto{PeanoInductionAddOne}[thm]{9,19}}
  \proofstepwidestarbreak[1,2,3,4,5,6,7]{a^{r+\lambda\cdot q}=a^r}{\rRE{\rUE{20},\rAEb{\rAEb{\rAEb{3}}}}}
\end{tabproofwide}


\hypertarget{index-period.proof.FiniteSemigroupIdempotentExponentData}{}
\IndexPeriodStatement{FiniteSemigroupIdempotentExponentData}
\begin{tabproofwide}
  \proofstepwidestarbreak[1]{\mu,\lambda\in\mathbb N}{\rA}
  \proofstepwidestarbreak[2]{\mu\neq0}{\rA}
  \proofstepwidestarbreak[3]{\lambda\neq0}{\rA}
  \proofstepwidestarbreak[1]{\lambda\cdot\mu\in\mathbb N}{\FormulaRefAuto{PeanoMulClosure}[thm]{\rAEb{1},\rAEa{1}}}
  \proofstepwidestarbreak[1--3]{\lambda\cdot\mu\neq0}{\FormulaRefAuto{PeanoMulNonzero}[thm]{\rAEb{1},\rAEa{1},3,2}}
  \proofstepwidestarbreak[1,3]{\exists k\in\mathbb N\,(\lambda=k+1)}{\FormulaRefAuto{PeanoNonzeroIsAddOne}[thm]{\rAEb{1},3}}
  \proofstepwidestarbreak[7]{k\in\mathbb N\land\lambda=k+1}{\rA}
  \proofstepwidestarbreak[1,7]{k\cdot\mu\in\mathbb N}{\FormulaRefAuto{PeanoMulClosure}[thm]{\rAEa{7},\rAEa{1}}}
  \proofstepwidestarbreak[1,7]{\lambda\cdot\mu=\mu+k\cdot\mu}{\rChain{\rIE{\rAEb{7}},\FormulaRefAuto{PeanoMulAddOneLeft}[thm]{\rAEa{7},\rAEa{1}},\FormulaRefAuto{PeanoAddCommutative}[thm]{8,\rAEa{1}}}}
  \proofstepwidestarbreak[1,7]{\exists t\in\mathbb N\,(\mu+t=\lambda\cdot\mu)}{\rEI{\rAI{8,\FormulaRefAuto{a=b\vdash b=a}[thm]{9}}}}
  \proofstepwidestarbreak[1,7]{\mu\leq\lambda\cdot\mu}{\FormulaRefAuto{NaturalOrderAdditiveCharacterization}[thm]{\rAEa{1},4,10}}
  \proofstepwidestarbreak[1,3]{\mu\leq\lambda\cdot\mu}{\rEE{6,7,11}}
  \proofstepwidestarbreak[1--3]{\lambda\cdot\mu\in\mathbb N\land\lambda\cdot\mu\neq0\land\mu\leq\lambda\cdot\mu}{\rAI{4,\rAI{5,12}}}
\end{tabproofwide}


\Needspace{16\baselineskip}
\hypertarget{index-period.proof.FiniteSemigroupIdempotentPower}{}
\IndexPeriodStatement{FiniteSemigroupIdempotentPower}
Innerhalb der Tabelle bedeutet \(n=\lambda\cdot\mu\).
\begin{tabproofwide}
  \proofstepwidestarbreak[1]{\FiniteSemigroupAlg{A}{\star}}{\rA}
  \proofstepwidestarbreak[2]{a\in A}{\rA}
  \proofstepwidestarbreak[1]{\SemigroupAlg{A}{\star}}{\FormulaRefAuto{FiniteSemigroupBaseAxiom}[ax]{1}}
  \proofstepwidestarbreak[1,2]{\exists\mu\,\exists\lambda\,\operatorname{IP}_{A,\star}(a;\mu,\lambda)}{\FormulaRefAuto{FiniteSemigroupIndexPeriodExistence}[thm]{1,2}}
  \proofstepwidestarbreak[5]{\exists\lambda\,\operatorname{IP}_{A,\star}(a;\mu,\lambda)}{\rA}
  \proofstepwidestarbreak[6]{\operatorname{IP}_{A,\star}(a;\mu,\lambda)}{\rA}
  \proofstepwidestarbreak[1,2,6]{\mu\in\mathbb N\land\mu\neq0\land\lambda\in\mathbb N\land\lambda\neq0}{\FormulaRefAuto{FiniteSemigroupIndexPeriodPositiveData}[ax]{3,2,6}}
  \proofstepwidestarbreak[1,2,6]{a^{\mu+\lambda}=a^\mu}{\FormulaRefAuto{FiniteSemigroupIndexPeriodEquality}[ax]{3,2,6}}
  \proofstepwidestarbreak[1,2,6]{n\in\mathbb N\land n\neq0\land\mu\leq n}{\FormulaRefAuto{FiniteSemigroupIdempotentExponentData}[thm]{\rAI{\rAEa{7},\rAEa{\rAEb{\rAEb{7}}}},\rAEa{\rAEb{7}},\rAEb{\rAEb{\rAEb{7}}}}}
  \proofstepwidestarbreak[1,2,6]{a^{n+n}=a^n}{\FormulaRefAuto{FiniteSemigroupMultiplePeriodShift}[thm]{3,2,\rAI{\rAEa{7},\rAI{\rAEa{\rAEb{\rAEb{7}}},\rAI{\rAEa{9},\rAEa{7}}}},\rAEa{\rAEb{7}},\rAEb{\rAEb{\rAEb{7}}},8,\rAEb{\rAEb{9}}}}
  \proofstepwidestarbreak[1,2,6]{a^{n+n}=a^n\star a^n}{\FormulaRefAuto{SemigroupPowerAdditionLaw}[thm]{3,2,\rAEa{9},\rAEa{9},\rAEa{\rAEb{9}},\rAEa{\rAEb{9}}}}
  \proofstepwidestarbreak[1,2,6]{a^n\in A}{\FormulaRefAuto{SemigroupPowerClosure}[thm]{3,2,\rAEa{9},\rAEa{\rAEb{9}}}}
  \proofstepwidestarbreak[1,2,6]{a^n\star a^n=a^n}{\rChain{\FormulaRefAuto{a=b\vdash b=a}[thm]{11},10}}
  \proofstepwidestarbreak[1,2,6]{a^n\in E_{A,\star}}{\FormulaRefAuto{SemigroupIdempotentMembership}[thm]{3,\rAI{12,13}}}
  \proofstepwidestarbreak[1,2,6]{\exists n\in\mathbb N\,(n\neq0\land a^n\in E_{A,\star})}{\rEI{\rAI{\rAEa{9},\rAI{\rAEa{\rAEb{9}},14}}}}
  \proofstepwidestarbreak[1,2,5]{\exists n\in\mathbb N\,(n\neq0\land a^n\in E_{A,\star})}{\rEE{5,6,15}}
  \proofstepwidestarbreak[1,2]{\exists n\in\mathbb N\,(n\neq0\land a^n\in E_{A,\star})}{\rEE{4,5,16}}
\end{tabproofwide}


\hypertarget{index-period.proof.FiniteSemigroupIdempotentExists}{}
\IndexPeriodStatement{FiniteSemigroupIdempotentExists}
% Aus Band 37 ausgelagerte Beweistabelle.
\begin{tabproofwide}
  \proofstepwidestar[1]{\FiniteSemigroupAlg{A}{\star}}{\rA}
  \proofstepwidestar[2]{A\neq\varnothing}{\rA}
  \proofstepwidestar[2]{\exists a\in A}{%
    \FormulaRefAuto{A\neq\varnothing\vdash\exists x\in A}[thm]{2}}
  \proofstepwidestar[4]{a\in A}{\rA}
  \proofstepwidestarbreak[1,4]{%
    \exists n\in\mathbb N\,
      \bigl(n\neq0\land a^n\in E_{A,\star}\bigr)}{%
    \FormulaRefAuto{FiniteSemigroupIdempotentPower}[thm]{1,4}}
  \proofstepwidestar[6]{%
    n\in\mathbb N\land n\neq0\land a^n\in E_{A,\star}}{\rA}
  \proofstepwidestar[6]{E_{A,\star}\neq\varnothing}{%
    \FormulaRefAuto{x\in A\vdash A\neq\varnothing}[thm]{%
      \rAEb{\rAEb{6}}}}
  \proofstepwidestar[1,4]{E_{A,\star}\neq\varnothing}{\rEE{5,6,7}}
  \proofstepwidestar[1,2]{E_{A,\star}\neq\varnothing}{\rEE{3,4,8}}
\end{tabproofwide}

